\section{Generators and the Cayley graph}\label{sec:cayley}

Three of the examples above were built by the same trick: take a geometric object, decorate it so
that unwanted symmetries are spoiled, and read off the group that survives. This section turns the
trick round. Instead of starting with an object and finding its group, start with a group and
\emph{build} an object.

\subsection{Generating sets}

\begin{defi}\label{def:gen}
A subset $S$ of a group $G$ \defn{generates} $G$ if every element of $G$ can be written as a product
of finitely many elements of $S$ and inverses of elements of $S$. Such a product is called a
\defn{word} in $S$, and $S$ is a \defn{generating set}.
\end{defi}

Every example so far came with one.

\begin{center}
\small
\begin{tabular}{@{}lll@{}}
\toprule
group & a generating set & typical element as a word\\
\midrule
$\Z$ & $\{1\}$ & $5=1+1+1+1+1$\\
$\Z/n\Z$ & $\{1\}$ & every element, by repeated addition\\
$\Z^2$ & $\{(1,0),(0,1)\}$ & $(2,1)=$ right, right, up\\
$D_n$ & $\{s,t\}$, two reflections & $(st)^m$ and $s(st)^m$\\
$S_n$ & the $n-1$ neighbour swaps & any permutation, by sorting\\
rotations of the cube & two quarter turns about different axes & every one of the $24$\\
$F_2$ & $\{a,b\}$ & the word itself\\
\bottomrule
\end{tabular}
\end{center}

A group can have many generating sets of very different sizes, and finding a small one is often the
first step in understanding it. Note that $S$ generating $G$ says nothing about \emph{uniqueness}:
in $\Z$ the element $0$ is both the empty product and $1+(-1)$.

\begin{tryit}
Show that $\{2,3\}$ generates $\Z$ but $\{2\}$ does not. Show that $\{4\}$ does not generate
$\Z/12\Z$ but $\{5\}$ does. Then convince yourself that the three face half turns of the cube
together with the identity form a group that no single element generates.
\end{tryit}

\subsection{The Cayley graph}

\begin{defi}\label{def:cayley}
Let $G$ be a group and $S$ a generating set. The \defn{Cayley graph} of $G$ with respect to $S$ is
the following decorated object. Put one dot for each element of $G$. For each element $g$ and each
generator $s\in S$, draw an arrow from the dot $g$ to the dot $gs$, coloured according to which
generator $s$ it came from.
\end{defi}

Three features of the definition deserve a sentence each. The dots are the group elements, so the
graph is as big as the group. The arrows are labelled by generators, so the same picture with a
different generating set looks different. And the arrows are coloured and directed, which is exactly
the decoration that spoils unwanted symmetries, as in Sections~\ref{sec:arrowgon} and \ref{sec:Z2}.

Every picture in the list below has already appeared in this chapter, without the name.

\begin{center}
\small
\begin{tabular}{@{}llll@{}}
\toprule
group & generating set & Cayley graph & drawn in\\
\midrule
$\Z/n\Z$ & $\{1\}$ & the arrowed $n$-gon & Figure~\ref{fig:arrowgon}\\
$\Z$ & $\{1\}$ & the arrowed line & Figure~\ref{fig:arrowline}\\
$\Z^2$ & $\{(1,0),(0,1)\}$ & the grid with coloured edges & Figure~\ref{fig:grid}\\
$F_2$ & $\{a,b\}$ & the $4$-valent tree & Figure~\ref{fig:tree}\\
$D_n$ & $\{s,t\}$ & a $2n$-gon with two edge colours & Problem~\ref{prob:cayleyD3}\\
\bottomrule
\end{tabular}
\end{center}

\begin{why}
\emph{Why the arrow goes from $g$ to $gs$ and not to $sg$.} Because of the convention of
Definition~\ref{def:order}. Walking along an $s$-arrow means multiplying on the \emph{right} by $s$.
Reading a path from the base dot $1$ then spells out a word in the generators from left to right,
which is exactly how words are written. If the arrows ran to $sg$ instead, reading a path would
spell words backwards.
\end{why}

\sfig{A group's multiplication table and its Cayley graph carry the same information. The table is
what you compute with; the graph is what you can look at.}{\figcayleytable}

\begin{bigpic}
\emph{The punchline of the office hour.} Take a group $G$, choose a generating set, and build the
Cayley graph. That graph is a decorated object, and one can ask for its symmetries: the motions of
the graph sending dots to dots, arrows to arrows, and respecting the colours. The answer is that
these symmetries are exactly the elements of $G$ again, with $g$ acting by sending each dot $x$ to
the dot $gx$. So \emph{every group is the symmetry group of some object}, and the object can be
manufactured from the group itself. The formal statement and proof belong to the next office hour;
what matters here is that the slogan ``groups are symmetries'' is a theorem, not a slogan.
\end{bigpic}

\begin{tryit}
Draw the Cayley graph of $\Z/6\Z$ with the generating set $\{1\}$, and then with $\{2,3\}$. The two
pictures look nothing alike, yet they describe the same group. Then draw the Cayley graph of $S_3$
with its two neighbour swaps; you should get a hexagon with two colours of edge.
\end{tryit}
